\documentclass[11pt]{article}
\usepackage[a4paper,margin=25mm]{geometry}
\usepackage[T1]{fontenc}
\usepackage[utf8]{inputenc}
\usepackage{amsmath,amssymb,booktabs}
\usepackage{enumitem}
\usepackage{array}
\usepackage{microtype}
\usepackage{xurl}
\usepackage[hidelinks]{hyperref}

\setlength{\parindent}{0pt}
\setlength{\parskip}{0.55em}
\setlist{nosep,leftmargin=*}
\emergencystretch=3em
\newcolumntype{L}[1]{>{\raggedright\arraybackslash}p{#1}}

\title{M13-C3 --- H1 support ablation\\
AAMAS nested-prefix missing \((0,0,0)\)}
\author{}
\date{H1 run / analysed / not a Bucket~3 claim}

\begin{document}
\maketitle

\section{Status}

\textbf{H1 run / analysed.}
Nested-prefix AAMAS ablation on fixture
\path{m13_bucket3_binary_collider_and} is complete.
Markdown twin: \texttt{h1\_support\_ablation.md}.
Closed H0 baseline preserved (\(n{=}30\), seed~42; see
\texttt{learning\_analysis.md}).
\textbf{Not a Bucket~3 claim.}
H2 is documented in \texttt{h2\_irrelevant\_covariate.md}; H3--H4 remain
deferred.

\section{Research question}

On the closed H0 AND-collider fixture under AAMAS, if a nested seed-42 prefix
omits the only sample atoms \((0,0,0)\), do root targets flip from the H0
outcome \texttt{completed\_no\_solution} to a \texttt{solved} delta that includes
a spurious both-0 co-occurrence Horn rule?

\section{Evaluative stance (do not invert)}

For this fixture, roots \(a\) and \(b\) have \textbf{no} deterministic causal
rule over the other observed variables.
The evaluator-only reference records
\texttt{no\_observed\_parent\_deterministic\_rule} for those targets.

\begin{center}
\begin{tabular}{@{}L{0.46\textwidth}L{0.44\textwidth}@{}}
\toprule
Outcome & Reading \\
\midrule
H0 AAMAS root \texttt{completed\_no\_solution} &
  \textbf{Desired / correct} for those targets \\
H1 AAMAS root \texttt{solved} with co-occurrence Horn rules &
  \textbf{Incorrect / spurious} --- a finite-sample mistake, \textbf{not}
  successful mechanism recovery \\
\bottomrule
\end{tabular}
\end{center}

H1 demonstrates that \textbf{AAMAS can incorrectly return a solution for a root
target when the sample undersamples rare joint assignments that would have
blocked an over-general rule}, especially when that assignment has \textbf{low
population probability}
(\(P(A{=}0,B{=}0,C{=}0)=3/50=0.06\)) and \textbf{\(n\) is small} (here
\(n=25\)).

Do \textbf{not} frame H1 as ``AAMAS learns roots better with less data''.
Do \textbf{not} treat root \texttt{solved} as progress.

\section{Design}

\subsection{Intervention}

Nested seed-42 prefix \(n{=}25\) of the closed H0 sample \texttt{n30\_seed42}
(exact first 25 rows).
That prefix omits the only H0 \((0,0,0)\) rows (CSV rows~\textbf{26} and
\textbf{30}).
Fixture, encoding, AAMAS configuration, and seed sequence are otherwise
unchanged.

Verified controls:
\begin{itemize}
  \item \texttt{n25\_seed42.csv} is a row-for-row prefix of
    \texttt{n30\_seed42.csv}
  \item No \((0,0,0)\) atom appears in \(n{=}25\) (joint counts:
    \((1,1,1):14\), \((1,0,0):7\), \((0,1,0):4\))
  \item Same AAMAS \texttt{configuration.id} (\texttt{aamas2025}) and
    configuration hash as H0
    (\texttt{sha256:93ee2bb829096988ba7b84a7ef2f0247e0cfd8a069909a5146500be9d10e7407}).
    Only the sample (\(n\), CSV hash) changes
  \item The population still assigns positive mass to \((0,0,0)\).
    Only the finite table omits it
\end{itemize}

\subsection{Paths}

\begin{itemize}
  \item Sample:
    \path{causal/outputs/causal_fixtures/m13_bucket3_binary_collider_and/samples/n25_seed42.csv}
    (+ \texttt{.manifest.json})
  \item Config:
    \path{causal/configs/targetwise/m13_bucket3_binary_collider_and/aamas2025/n25_seed42.yaml}
  \item H1 collection:
    \path{causal/outputs/aba_learning/targetwise/m13_bucket3_binary_collider_and/aamas2025/n25_seed42/}
  \item H0 AAMAS comparison (unchanged):
    \path{.../aamas2025/n30_seed42/}
\end{itemize}

Sample hash (H1):
\texttt{sha256:c7ccf94c10a7862f34b7eeedb61529cf60ac0b13b932ac608beaea1a6d17fdd4}.

\subsection{Command (reproducible form)}
\begin{verbatim}
python -m causal.targetwise.cli run \
  --config causal/configs/targetwise/m13_bucket3_binary_collider_and/aamas2025/n25_seed42.yaml
\end{verbatim}
Learner options remain those of \texttt{configs/aamas2025\_config.pl}: brave,
greedy, mgr, \texttt{folding\_space(bk)}, relto, \texttt{check\_ic}.

\section{Results (AAMAS)}

\begin{center}
\begin{tabular}{@{}llll@{}}
\toprule
Target & H0 \(n{=}30\) & H1 \(n{=}25\) & H1 reading \\
\midrule
\(a\) & no-solution (\textbf{correct}) &
  \texttt{solved} (\textbf{incorrect}) & spurious root solve \\
\(b\) & no-solution (\textbf{correct}) &
  \texttt{solved} (\textbf{incorrect}) & spurious root solve \\
\(c\) & solved AND (control) & solved AND & control unchanged in kind \\
\bottomrule
\end{tabular}
\end{center}

H1 \(E^\pm\): \(a\) 21/4, \(b\) 18/7, \(c\) 14/11.
Artefact audit SAT. Zero assumptions on all three H1 deltas.

Root \(\Delta\) for \(a\):
\begin{verbatim}
a(A) :- b_val_1(A), c_val_1(A).
a(A) :- b_val_0(A), c_val_0(A).
\end{verbatim}
Target \(b\) is the mirror with \(a\)/\(c\).
Target \(c\):
\begin{verbatim}
c(A) :- a_val_1(A), b_val_1(A).
\end{verbatim}
The both-0 root rules are \textbf{false} as population implications.
They are sample-adequate only because the H1 table contains no opposite-label
witness in that predictor cell.

\section{Trace / fixture reading}

Procedural comparison uses
\path{.../aamas2025/n25_seed42/cells/target-a/output/prolog.stdout}
against the closed H0 trace under \texttt{n30\_seed42}
(target \(b\) is symmetric).

Under AAMAS greedy folding, both runs build maximal co-occurrence bodies from
BK.
The both-1 body
\texttt{a(A) :- b\_val\_1(A), c\_val\_1(A).}
is accepted under brave entailment on both tables.

The both-0 body
\texttt{a(A) :- b\_val\_0(A), c\_val\_0(A).}
is then proposed from an early both-0 positive (row~8 on both traces).

\begin{itemize}
  \item On \textbf{H1} (\(n{=}25\)): brave entailment of
    \(\langle E^+,E^-\rangle\) \textbf{succeeds} immediately after that fold.
    No assumption is introduced.
    The queue empties.
    The run terminates as \texttt{solved} with the two Horn rules above.
  \item On \textbf{H0} (\(n{=}30\)): the same both-0 body fails brave
    entailment because rows~\textbf{26} and~\textbf{30} are \(E^-\) examples
    that share \(\mathtt{b\_val\_0}\) and \(\mathtt{c\_val\_0}\).
    Assumption/contrary repair reconstructs the same both-0 body for the
    contrary.
    Relto cannot introduce a further assumption.
    The engine reports KO and \texttt{* No solution found!}.
\end{itemize}

The AAMAS search procedure is therefore the same until the both-0 commit.
What flips the outcome is the \textbf{absence versus presence of a label
conflict inside that greedy candidate body cell}.
The missing conflict is exactly the rare joint \((0,0,0)\), which has
population probability \(3/50=0.06\) and appears only twice in the H0 table of
size~30.

Child target \(c\) remains a control: both H0 and H1 emit the assumption-free
AND conjunction.
That does not change the root evaluative stance.

\section{Bounded interpretation}

On this fixture and AAMAS configuration, omitting the rare conflicting joint
from a small nested prefix converts the \textbf{correct} root outcome
(\texttt{completed\_no\_solution}) into an \textbf{incorrect} spurious Horn
covering of the ambiguous both-0 cell.
That is a finite-sample / support-sensitivity effect for greedy maximal
co-occurrence search.
It is \textbf{not} evidence that shorter tables recover root mechanisms.

Roots remain non-mechanistic under the evaluator reference.
Predictive co-occurrence covering is not mechanism recovery.

\section{Explicit non-claims}

\begin{itemize}
  \item Not a Bucket~3 claim.
  \item Not Russo-style Causal ABA, DAG recovery, or CPDAG recovery.
  \item Not ``AAMAS recovered the root mechanism'' on \(n{=}25\).
  \item Not ``AAMAS learns roots better with less data''.
  \item Not a claim that every seed or every \(n\) behaves identically.
  \item Not a claim that omitting \((0,0,0)\) is the only way AAMAS can emit a
    spurious root rule.
  \item H3--H4 are untouched by this record (H2 is a separate later probe).
\end{itemize}

\section{Cross-links}

\begin{itemize}
  \item Investigation hub: \texttt{experiment.md}
  \item Closed H0 learning record: \texttt{learning\_analysis.md} / \texttt{.tex}
  \item Probe catalogue: \texttt{future\_probes.md} / \texttt{.tex}
  \item Bucket~3 planning hub: \texttt{../M1.3-bucket3-claims.md}
    (still \textbf{no claim})
\end{itemize}

\section{Next}

H1 documentation is complete.
H1 documentation is complete.
Probe H2 is documented in \texttt{h2\_irrelevant\_covariate.md}.
Remaining probes H3--H4 stay deferred as recorded in
\texttt{future\_probes.md}.
Still \textbf{no Bucket~3 claim}.

\end{document}
